\documentclass[10pt,a4paper]{amsart}
\usepackage[margin=19mm]{geometry}
\usepackage{amsmath,amssymb,mathtools}
\usepackage[scr=rsfs,cal=euler]{mathalfa}
\usepackage{xfrac}
\usepackage[unicode,hidelinks,bookmarks=false]{hyperref}
\newcommand{\ie}{i.e.\ }
\newcommand{\cf}{cf.\ }
\newcommand{\resp}{respectively\ }
\newcommand{\Subsection}{\subsection}
\providecommand{\nobreakdash}{\nobreak\hbox{-}}
\setlength{\emergencystretch}{2em}
\multlinegap0pt
\usepackage{amssymb}
\newtheorem{lem}{Lemma}
\def \Z {\mathbb Z}
\def \d {\delta}
\def \ep {\epsilon}
\title{Realizing Arbitrary Depth}
\author{Ethan Akin}
\date{Source: arXiv:2606.01499v2 (CC BY 4.0)}
\begin{document}
\maketitle

\noindent\emph{Source and licence.} Realizing Arbitrary Depth. Version arXiv:2606.01499v2 (CC BY 4.0), distributed by arXiv under the Creative Commons Attribution 4.0 International licence, http://creativecommons.org/licenses/by/4.0/, as recorded in the arXiv licence metadata for this exact version and shown on the abstract page https://arxiv.org/abs/2606.01499v2. Full licence notice: \texttt{licenses/arxiv-2606.01499-CC-BY-4.0.txt}.\par\smallskip

\noindent\textit{Editorial scope.} Counted unit: the article's lem lem1.04 (source0.tex lines 370--370). The article's complete original proof of it is delivered with it (source0.tex lines 372--383, one \texttt{proof} environment of 12 lines). No mathematics is written, repaired or re-derived: the statement and the proof are the article's own lines, in the article's own notation, and no retained line is substituted or re-flowed.  No further source-local context is needed: every cross-reference of every retained line resolves inside the two delivered spans.  Nothing was omitted from any retained span, so the package carries no omission record.  No retained line contains a citation, so the wrapper carries no bibliography and no bibliography entry.  No retained line contains a figure environment, an image-inclusion command or drawing code, so no image is delivered and the package declares no asset dependency.  Editorial formatting only, in addition: an AMS \texttt{amsart} preamble that restates the article's own declarations and macros used by the retained lines, namely the environments \texttt{lem} on separate counters, together with the standard packages those lines need.  The retained environments print in this document as Lemma 1 in order of appearance, whereas the article prints the same items with the numbering of its own full document.  The complete native source of the article as distributed in the arXiv e-print for this exact version is bundled unchanged as \texttt{sources/201833/source0.tex}.
\begin{lem}\label{lem1.04} A dynamical system $(X,f)$ is chain transitive if and only if it admits a full, asymptotic chain. \end{lem}

\begin{proof} If $\{ x_k : k \in \Z \}$ is a full asymptotic chain, $y_1, y_2 \in X$, and $\d < \ep/2$ an $\ep/2$ modulus of uniform
continuity for $f$, then we can choose $N$ so that  $d(f(x_{k-1}),x_k) < \d$ for all $k > N$. By the density condition, there exist $k_1, k_2$ with
$N \le k_1  < k_2 - 2$ and such that $d(x_{k_1},y_1) < \d$ and $d(x_{k_2},y_2) < \d$. Then $\{ y_1, x_{k_1+1}, \dots x_{k_2 - 1}, y_2 \}$
is an $\ep$ chain from $y_1$ to $y_2$.

A compact metric space is totally bounded and so there exists for every $\ep > 0$ a finite sequence in $X$ which is $\ep$ dense,
i.e. it passes within $\ep$ of every point of $X$.
By concatenating such finite sequences it is easy to obtain a bi-infinite sequence $\{ y_k : k \in \Z \}$ with dense tails. If $(X,f)$ is chain
transitive, then we can insert between $y_{k}$ and $y_{k+1}$ a $1/(|k| + 1)$ chain from $y_{k}$ and $y_{k+1}$. After renumbering we obtain
a full asymptotic chain.

\end{proof}\vspace{1cm}

\end{document}
